% GENERATED FILE. Edit canonical lessons, then run npm run generate:books.
% canonical-source-hash: fnv1a64:da262558218a2ee2
% canonical-lessons: LA-C253-feriae, LA-C253-convivium, LA-C253-sacrificium, LA-C253-ara, LA-C253-dea

\chapter{Holidays, Banquet, Sacrifice, Altar, Goddess}
\label{ch:la-holidays-banquet-sacrifice-altar-goddess}
\hlchaptermodality{253}

\begin{chapteropening}
\textbf{By the end of this chapter:} \emph{I can say holidays, a banquet, a sacrifice, an altar and a goddess.}

Complete the last lesson of chapter 253: I can say holidays, a banquet, a sacrifice, an altar and a goddess.
\end{chapteropening}

\section[fēriae, fēriārum]{fēriae, fēriārum — holidays}

\label{lesson:LA-C253-feriae}



\begin{quote}
Before the new one: say the Latin for a candlestick, then the Latin for a small tablet.
\end{quote}

\subsection*{You'll want to know: fēriae, fēriārum}
\textbf{fēriae, fēriārum} — \textquotedblleft{}holidays\textquotedblright{}.

\begin{cousinweb}
\textbf{fēriae} comes from an older \emph{fēsiae} (Lewis \& Short), from the same root as \emph{fēstus}, \textquotedblleft{}festive\textquotedblright{}. Both are usually traced to Proto-Indo-European \emph{*dheh1s-}, \textquotedblleft{}god, sacred\textquotedblright{} (EtymDB 2.1). English \emph{fair}, a market, comes from it through French.
\end{cousinweb}

\begin{grammarlens}[title={What you've built}]
One more word for this chapter.
\end{grammarlens}

\subsection*{Your turn}
\begin{itemize}
\raggedright
  \item \emph{Say it:} \emph{fēriae, fēriārum}
  \item \emph{Say it:} \emph{fēriae, fēriārum}, once more
  \item \emph{Say it:} say \emph{fēriae, fēriārum}
  \item \emph{Recall:} say the Latin for a candlestick, then the Latin for a small tablet, then say \emph{fēriae, fēriārum} again
\end{itemize}

\begin{tcolorbox}[breakable,colback=teal!4,colframe=teal!35!black,title={Before you move on}]
What does fēriae, fēriārum mean? (\textquotedblleft{}Holidays\textquotedblright{}.) Say it once more.
\end{tcolorbox}
\section[convīvium, convīviī]{convīvium, convīviī — a banquet}

\label{lesson:LA-C253-convivium}



\begin{quote}
Before the new one: say the Latin for a small tablet, then the Latin for holidays.
\end{quote}

\subsection*{You'll want to know: convīvium, convīviī}
\textbf{convīvium, convīviī} — \textquotedblleft{}a banquet\textquotedblright{}.

\begin{cousinweb}
\textbf{convīvium} is \emph{con-}, \textquotedblleft{}together\textquotedblright{}, with the root of \emph{vīvō}, \textquotedblleft{}to live\textquotedblright{} (Lewis \& Short): living together, then a shared meal.
\end{cousinweb}

\begin{grammarlens}[title={What you've built}]
One more word for this chapter.
\end{grammarlens}

\subsection*{Your turn}
\begin{itemize}
\raggedright
  \item \emph{Say it:} \emph{convīvium, convīviī}
  \item \emph{Say it:} \emph{convīvium, convīviī}, once more
  \item \emph{Say it:} say \emph{convīvium, convīviī}
  \item \emph{Recall:} say the Latin for a small tablet, then the Latin for holidays, then say \emph{convīvium, convīviī} again
\end{itemize}

\begin{tcolorbox}[breakable,colback=teal!4,colframe=teal!35!black,title={Before you move on}]
What does convīvium, convīviī mean? (\textquotedblleft{}A banquet\textquotedblright{}.) Say it once more.
\end{tcolorbox}
\section[sacrificium, sacrificiī]{sacrificium, sacrificiī — a sacrifice}

\label{lesson:LA-C253-sacrificium}



\begin{quote}
Before the new one: say the Latin for holidays, then the Latin for a banquet.
\end{quote}

\subsection*{You'll want to know: sacrificium, sacrificiī}
\textbf{sacrificium, sacrificiī} — \textquotedblleft{}a sacrifice\textquotedblright{}.

\begin{cousinweb}
\textbf{sacrificium} is built from \emph{sacer}, \textquotedblleft{}holy\textquotedblright{}, and \emph{faciō}, \textquotedblleft{}to make\textquotedblright{} (EtymDB 2.1): to make something holy. \emph{sacer} is usually traced to Proto-Indo-European \emph{*sh2kros}, \textquotedblleft{}sacred\textquotedblright{} (EtymDB 2.1).
\end{cousinweb}

\begin{grammarlens}[title={What you've built}]
One more word for this chapter.
\end{grammarlens}

\subsection*{Your turn}
\begin{itemize}
\raggedright
  \item \emph{Say it:} \emph{sacrificium, sacrificiī}
  \item \emph{Say it:} \emph{sacrificium, sacrificiī}, once more
  \item \emph{Say it:} say \emph{sacrificium, sacrificiī}
  \item \emph{Recall:} say the Latin for holidays, then the Latin for a banquet, then say \emph{sacrificium, sacrificiī} again
\end{itemize}

\begin{tcolorbox}[breakable,colback=teal!4,colframe=teal!35!black,title={Before you move on}]
What does sacrificium, sacrificiī mean? (\textquotedblleft{}A sacrifice\textquotedblright{}.) Say it once more.
\end{tcolorbox}
\section[āra, ārae]{āra, ārae — an altar}

\label{lesson:LA-C253-ara}



\begin{quote}
Before the new one: say the Latin for a banquet, then the Latin for a sacrifice.
\end{quote}

\subsection*{You'll want to know: āra, ārae}
\textbf{āra, ārae} — \textquotedblleft{}an altar\textquotedblright{}.

\begin{cousinweb}
\textbf{āra} comes from an older \emph{āsa}, as Oscan \emph{aasa}, \textquotedblleft{}altar\textquotedblright{}, shows (Lewis \& Short); it is usually traced to Proto-Indo-European \emph{*h2eHs-}, \textquotedblleft{}to burn, a hearth\textquotedblright{} (EtymDB 2.1).
\end{cousinweb}

\begin{grammarlens}[title={What you've built}]
One more word for this chapter.
\end{grammarlens}

\subsection*{Your turn}
\begin{itemize}
\raggedright
  \item \emph{Say it:} \emph{āra, ārae}
  \item \emph{Say it:} \emph{āra, ārae}, once more
  \item \emph{Say it:} say \emph{āra, ārae}
  \item \emph{Recall:} say the Latin for a banquet, then the Latin for a sacrifice, then say \emph{āra, ārae} again
\end{itemize}

\begin{tcolorbox}[breakable,colback=teal!4,colframe=teal!35!black,title={Before you move on}]
What does āra, ārae mean? (\textquotedblleft{}An altar\textquotedblright{}.) Say it once more.
\end{tcolorbox}
\section[dea, deae]{dea, deae — a goddess}

\label{lesson:LA-C253-dea}



\begin{quote}
Before the new one: say the Latin for a sacrifice, then the Latin for an altar.
\end{quote}

\subsection*{You'll want to know: dea, deae}
\textbf{dea, deae} — \textquotedblleft{}a goddess\textquotedblright{}.

\begin{cousinweb}
\textbf{dea} is the feminine of \emph{deus}, \textquotedblleft{}a god\textquotedblright{} (Lewis \& Short); Old Latin wrote it \emph{deiva} (EtymDB 2.1).
\end{cousinweb}

\begin{grammarlens}[title={What you've built}]
One more word for this chapter.
\end{grammarlens}

\subsection*{Your turn}
\begin{itemize}
\raggedright
  \item \emph{Say it:} \emph{dea, deae}
  \item \emph{Say it:} \emph{dea, deae}, once more
  \item \emph{Say it:} say \emph{dea, deae}
  \item \emph{Recall:} say the Latin for a sacrifice, then the Latin for an altar, then say \emph{dea, deae} again
\end{itemize}

\begin{tcolorbox}[breakable,colback=teal!4,colframe=teal!35!black,title={Before you move on}]
What does dea, deae mean? (\textquotedblleft{}A goddess\textquotedblright{}.) Say it once more.
\end{tcolorbox}
